\documentclass[11pt]{article}
\usepackage[margin=0.9in]{geometry}
\usepackage{booktabs}
\usepackage[hidelinks]{hyperref}
\usepackage{parskip}
\usepackage{microtype}

\newcommand{\repourl}{\url{https://github.com/Jotsna22/cyberrisk-llm}}

\title{}
\author{}
\date{}

\begin{document}

% ---------------- Cover page ----------------
\thispagestyle{empty}
\begin{center}
\vspace*{2.5cm}
{\LARGE\bfseries CyberRisk-LLM}\\[0.4cm]
{\Large LLM-Assisted Attack Path Risk Assessment and Adaptive Defense}\\[1.2cm]
{\large CSCE 703 Project --- Final Report (Software/Tool Development)}\\[2cm]
{\large Team Members}\\[0.3cm]
{\normalsize
Srija Pentyala\\
Siva Sai Deepank Manoj\\
Ubaid Khan Mohammed\\
Shivani Suresh\\
Lakshmi Jotsna Madhavarapu\\
Vyshnavi Pamidi}\\[1.5cm]
{\normalsize Source code: \repourl}\\[0.3cm]
{\normalsize November 2026}
\end{center}

\newpage
\setcounter{page}{1}

% ---------------- Body ----------------
\begin{center}
{\Large\bfseries CyberRisk-LLM: LLM-Assisted Attack Path Risk Assessment and Adaptive Defense}
\end{center}

\section{Problem}
Security teams hold vulnerability descriptions, asset inventories, access rules, and policies in separate places. CVSS scores rate each weakness on its own, and do not show how an attacker chains weaknesses across hosts, which critical asset is reachable, or which control to apply first. Our project asks whether an LLM can convert such unstructured text into a structured model that supports attack-path discovery, risk ranking, and defense selection, and whether the downstream reasoning can stay deterministic and testable.

\section{Methodology and System Design}
The system has five stages. An extractor turns analyst text into an \emph{environment} model: assets with kind, criticality, and exposure; the vulnerabilities on each asset; and the directed links the firewall permits. An \emph{attack graph} adds an Internet node and an edge $u \rightarrow v$ wherever $u$ can reach $v$ and $v$ has a vulnerability. Each edge has a probability equal to the best exploitability of $v$'s weaknesses, where exploitability is the product of an attack-vector factor (network 0.9, adjacent 0.5, local 0.1) and a privilege factor (none 1.0, low 0.7, high 0.4). A path's probability is the product of its edge probabilities, and its \emph{risk} is that probability multiplied by the target's criticality and its worst CVSS impact. Total risk sums, over critical assets, the highest path risk reaching each one. Paths of any length are enumerated. An earlier version capped paths at four edges, which silently removed two critical assets five hops from the Internet (\texttt{VAULT-01} in the cloud template, \texttt{RES-DATA} in the distributed one) and understated risk in 46 of 80 environments. All reported results are without the cap. A \emph{defense planner} chooses patches (cost 1) and link blocks (cost 2) greedily, by risk reduction per unit cost, under a budget of 4. Language understanding is isolated from calculation: the LLM only produces the model, and every later step is deterministic.

The risk score ranks the worst path to each asset. It is not a probability of compromise, because a second route to the same asset adds nothing to it. We therefore also run an \emph{attacker simulation} on the ground truth. An attacker starts on the Internet, and every attack-graph edge is one exploitation attempt that succeeds independently with the edge's probability. Every asset reachable over successful edges is compromised, by any route and over any path length. We report \emph{attack success} (the probability that at least one critical asset is compromised), the expected number of \emph{critical assets compromised}, and \emph{expected loss} (the expected sum of criticality $\times$ CVSS impact over compromised critical assets). Our graphs have at most 11 edges, so all $2^{11}$ outcomes are enumerated exactly. Larger graphs fall back to seeded Monte Carlo.

\section{Technical Implementation}
The implementation is in Python (3.10 or later) and uses NetworkX for graph enumeration. It is organized into a data model (\texttt{models.py}), graph and risk engine (\texttt{graph.py}), planner (\texttt{defense.py}), baselines (\texttt{baselines.py}), a synthetic generator and report renderer (\texttt{synthetic.py}), extractors (\texttt{extraction.py}), metrics, an attacker simulation (\texttt{simulation.py}), statistics (\texttt{stats.py}), and an experiment runner. Two extractors are provided. The \emph{rule-based} extractor parses canonical phrasings with regular expressions, requires no model, and serves as the no-LLM baseline. The \emph{template parser} extends it to all twelve phrasings the easy report generator uses, and is the strongest hand-written parser for that format. The \emph{LLM extractor} requests structured JSON from a language model, validates every field, counts rejected facts, and passes the analyst text as data, not as instructions. Unusable output (invalid JSON, or an empty response) is treated as an empty model and recorded as an extraction failure, so one bad response cannot abort a run. Before planning, a set of sanity checks inspects the model. They flag a model with no internet-exposed entry point, no high-criticality asset, or links to assets missing from the inventory, because in each case the planner would otherwise recommend nothing without any warning. Gemini responses are cached on disk, keyed by model and report text, and the cache is committed, so the LLM results reproduce offline without an API key. The CLI offers \texttt{cyberrisk demo} (with \texttt{--extractor oracle}, \texttt{rules}, or \texttt{gemini}) and \texttt{cyberrisk evaluate}. A suite of 105 pytest tests covers risk arithmetic (including five-hop paths), the hard-report generator and trap scoring, extraction (including malformed model output), planner budgets, the sanity checks, the attacker simulation (hand-calculated cases, and exact against Monte Carlo), the statistics (against SciPy reference values), deterministic generation, metrics, the CLI, and offline reproduction from the cache.

\section{Experimental Setup}
We generated 80 environments: four topology templates (a small research lab, an HPC cluster, a cloud deployment, and a distributed multi-site network), each with 20 seeds. All assets, CVE identifiers (\texttt{SYN-CVE-xxxx}), and scores are fictional. We compared four configurations. \texttt{cvss\_only} patches the highest-CVSS vulnerabilities first. \texttt{rule\_based} patches exposed-asset vulnerabilities with CVSS $\geq$ 7 first, then the rest by CVSS. Both baselines use the true vulnerability list and exposure flags, which is a favorable assumption for them. \texttt{graph\_rule\_extraction} plans on the rule-based extraction of the report text. \texttt{graph\_oracle} plans on the true model and serves as an upper bound. \texttt{graph\_llm\_extraction} plans on the model that \texttt{gemini-3.5-flash-lite} extracts from the same report text (temperature 0). Every plan is applied to the ground-truth environment, so extraction errors appear as wasted or missing controls. Risk reduction is the fractional drop in total risk. Besides set F1 over extracted facts, we score \emph{attribute accuracy} on the facts both models share: criticality for assets, and CVSS, attack vector, and privilege for vulnerabilities. Set F1 alone cannot see a criticality downgrade.

\emph{Hard reports.} The easy reports write each fact type in only three fixed phrasings, so the template parser reads them perfectly, and they cannot show that an LLM is needed. We therefore render each of the 80 environments a second time as a \emph{hard} report. These use free-form prose with synonyms, markdown tables for inventory, scanner output and firewall rules, several facts per sentence, direction reversals, negations and deny rules. Four statements are contradicted later in the text: a CVSS rescored after vendor review, a firewall rule removed, a false positive withdrawn, and a decommissioned host. Every statement that disagrees with the ground truth is retracted later or is a planted injection, so the correct reading of each report is exactly the true environment. A \emph{hard + injections} variant adds three injections placed mid-report. The first is an instruction addressed to automated tools to record a critical asset as low criticality, which the extraction prompt says to ignore. The second is a forwarded message, explicitly marked unverified, claiming such a downgrade. The third is a change-log line, worded like the genuine corrections, saying an internet-facing host no longer accepts internet traffic. Each trap is scored on its own as correct, fooled, or untested. Gemini read all 160 hard reports (cached, temperature 0). We also sweep the budget from 1 to 4 for the baselines and the graph planner. All configurations run on the same environments, so we compare them with paired tests on the per-environment difference: 95\% percentile bootstrap intervals (5{,}000 resamples, resampling seeds within each topology template) and two-sided Wilcoxon signed-rank tests (normal approximation with tie correction).

\section{Results}
Table~\ref{tab:defense} reports defense outcomes, and Table~\ref{tab:sig} tests the differences. Table~\ref{tab:hard} reports the hard reports. Table~\ref{tab:sweep} varies the budget. Table~\ref{tab:extract} reports extraction quality. Table~\ref{tab:robust} perturbs the planner's input model directly.

\begin{table}[h]
\centering
\small
\setlength{\tabcolsep}{4pt}
\begin{tabular}{lccc}
\toprule
Configuration & Risk reduction [95\% CI] & Attack success & Loss reduction [95\% CI] \\
\midrule
\texttt{cvss\_only} & $0.771$ [0.715, 0.824] & $0.509 \rightarrow 0.177$ & $0.782$ [0.733, 0.829] \\
\texttt{rule\_based} & $0.894$ [0.858, 0.927] & $0.509 \rightarrow 0.051$ & $0.902$ [0.869, 0.932] \\
\texttt{graph\_rule\_extraction} & $0.000$ & $0.509 \rightarrow 0.509$ & $0.000$ \\
\texttt{graph\_template\_extraction} & $0.999$ [0.996, 1.000] & $0.509 \rightarrow 0.004$ & $0.998$ [0.995, 1.000] \\
\texttt{graph\_oracle} (upper bound) & $0.999$ [0.996, 1.000] & $0.509 \rightarrow 0.004$ & $0.998$ [0.995, 1.000] \\
\texttt{graph\_llm\_extraction} (Gemini) & $0.999$ [0.996, 1.000] & $0.509 \rightarrow 0.004$ & $0.998$ [0.995, 1.000] \\
\bottomrule
\end{tabular}
\caption{Defense outcomes over 80 synthetic environments, budget 4. Attack success (before $\rightarrow$ after the plan) and expected-loss reduction come from the attacker simulation on the ground truth. The expected number of critical assets compromised falls from 1.008 to 0.264 (\texttt{cvss\_only}), 0.083 (\texttt{rule\_based}) and 0.004 (graph planner). The Gemini plans are identical to the oracle plans in all 80 environments.}
\label{tab:defense}
\end{table}

\begin{table}[h]
\centering
\small
\begin{tabular}{llccc}
\toprule
Budget & Comparison (A $-$ B) & Mean difference [95\% CI] & A better / tie / B better & Wilcoxon $p$ \\
\midrule
4 & graph $-$ \texttt{rule\_based} & $+0.105$ [$+0.071$, $+0.141$] & 42 / 36 / 2 & $2.3 \times 10^{-8}$ \\
4 & graph $-$ \texttt{cvss\_only} & $+0.227$ [$+0.174$, $+0.284$] & 61 / 19 / 0 & $1.1 \times 10^{-11}$ \\
4 & \texttt{rule\_based} $-$ \texttt{cvss\_only} & $+0.123$ [$+0.065$, $+0.182$] & 43 / 23 / 14 & $3.3 \times 10^{-4}$ \\
1 & graph $-$ \texttt{rule\_based} & $+0.471$ [$+0.401$, $+0.539$] & 67 / 13 / 0 & $1.1 \times 10^{-12}$ \\
2 & graph $-$ \texttt{rule\_based} & $+0.264$ [$+0.197$, $+0.333$] & 49 / 18 / 13 & $2.6 \times 10^{-7}$ \\
\bottomrule
\end{tabular}
\caption{Paired comparisons of risk reduction on the same 80 environments (graph = planner on the true model). Expected-loss reduction gives the same conclusions (e.g.\ $+0.096$ [$+0.066$, $+0.130$] against \texttt{rule\_based} at budget 4). Seeds within a template share a network layout, so these tests describe the four topologies studied, not networks in general. Full results: \texttt{results/significance.csv}.}
\label{tab:sig}
\end{table}

\begin{table}[h]
\centering
\small
\setlength{\tabcolsep}{4pt}
\begin{tabular}{llccccc}
\toprule
Report & Extractor & Vuln.\ F1 & Link F1 & Exp.\ F1 & Risk reduction [95\% CI] & Attack succ. \\
\midrule
easy & template parser & 1.000 & 1.000 & 1.000 & 0.999 [0.996, 1.000] & 0.004 \\
easy & Gemini & 0.9995 & 0.998 & 1.000 & 0.999 [0.996, 1.000] & 0.004 \\
hard & template parser & 0.000 & 0.000 & 0.000 & 0.000 & 0.509 \\
hard & Gemini & 0.987 & 0.934 & 0.988 & 0.966 [0.928, 0.995] & 0.017 \\
hard + inj. & template parser & 0.000 & 0.000 & 0.000 & 0.000 & 0.509 \\
hard + inj. & Gemini & 0.999 & 0.955 & 0.204 & 0.195 [0.129, 0.266] & 0.436 \\
\bottomrule
\end{tabular}
\caption{Extraction and defense outcomes on the easy and hard reports (80 environments each). ``Attack succ.'' is attack success after the plan; ``hard + inj.'' is the hard report with injections. The rule-based extractor also scores 0 on both hard variants. Gemini was fooled by: a deny rule in 32\% of hard reports, a removed firewall rule in 20\%, and no rescored CVSS, withdrawn finding, negation or decommissioned host. With injections, it followed the instruction to automated tools in 40\% of reports, the unverified message in 11\%, and the plausible change-log update in 96\%. Per-trap counts: \texttt{results/hard\_text\_traps.csv}.}
\label{tab:hard}
\end{table}

\begin{table}[h]
\centering
\begin{tabular}{lccc}
\toprule
Budget & \texttt{cvss\_only} & \texttt{rule\_based} & graph planner (true model) \\
\midrule
1 & $0.294 \pm 0.295$ / 0.410 & $0.309 \pm 0.331$ / 0.337 & $0.780 \pm 0.215$ / 0.168 \\
2 & $0.506 \pm 0.316$ / 0.328 & $0.690 \pm 0.354$ / 0.143 & $0.954 \pm 0.095$ / 0.056 \\
3 & $0.636 \pm 0.327$ / 0.256 & $0.823 \pm 0.272$ / 0.079 & $0.993 \pm 0.017$ / 0.013 \\
4 & $0.771 \pm 0.276$ / 0.177 & $0.894 \pm 0.183$ / 0.051 & $0.999 \pm 0.009$ / 0.004 \\
\bottomrule
\end{tabular}
\caption{Risk reduction (mean $\pm$ sd) / attack success after the plan, as the budget grows (80 environments). The advantage of topology-aware planning is largest when few controls can be afforded.}
\label{tab:sweep}
\end{table}

\begin{table}[h]
\centering
\begin{tabular}{lcccccc}
\toprule
Extractor & Assets F1 & Exposure F1 & Vuln.\ F1 & Link F1 & Crit.\ acc. & Flagged \\
\midrule
Rule-based & 0.822 & 0.329 & 0.493 & 0.514 & 0.610 & 80/80 \\
Gemini & 1.000 & 1.000 & 0.9995 & 0.998 & 1.000 & 0/80 \\
\bottomrule
\end{tabular}
\caption{Extraction quality. ``Crit.\ acc.'' is criticality accuracy on assets both models contain. ``Flagged'' counts models that the sanity checks warn about. Vulnerability attribute accuracy (CVSS, vector, privilege) is 1.000 for both. Roughly two thirds of facts in the test text are paraphrased in ways the rule set does not match, so rule-based recall is low by design.}
\label{tab:extract}
\end{table}

\begin{table}[h]
\centering
\begin{tabular}{lcc}
\toprule
Perturbation of planner input & Risk reduction (mean $\pm$ sd) & Controls used \\
\midrule
clean & $0.999 \pm 0.009$ & 2.05 \\
25\% of links and vulnerabilities missing & $0.322 \pm 0.416$ & 0.93 \\
50\% of links and vulnerabilities missing & $0.029 \pm 0.156$ & 0.14 \\
Two invented links, one invented vulnerability & $0.991 \pm 0.040$ & 2.48 \\
Critical assets relabeled low (adversarial) & $0.000 \pm 0.000$ & 0.00 \\
\bottomrule
\end{tabular}
\caption{Model-level robustness of the graph planner (perfect model otherwise). The sanity checks flag all 80 adversarial models and none of the clean ones.}
\label{tab:robust}
\end{table}

\paragraph{LLM extraction (live run).} We ran \texttt{gemini-3.5-flash-lite} on the same 80 reports (temperature 0, cached responses in \texttt{results/llm\_cache/}). Extraction F1 averaged 1.00 for assets and exposure, 0.9995 for vulnerabilities (minimum 0.96), and 0.998 for links (minimum 0.92). Criticality and vulnerability attributes were all correct. Plans built on these models matched the oracle ($0.999 \pm 0.009$). In one response (small lab, seed 15) the model returned a link with a null endpoint. The sanity checks flagged it, and the parser now rejects such links instead of creating an asset named ``None''. The reports are templated synthetic text, so this shows the model can read the format, not that it would handle real analyst reports. On a smaller robustness run (4 templates, 5 seeds each; \texttt{results/llm\_robustness.csv}) the LLM behaved as follows: with 30\% of structural lines removed, vulnerability and link F1 fell to 0.83 and 0.82 and risk reduction fell to 0.18; with three invented lines appended, asset F1 fell to 0.91 (the phantom host was extracted) but risk reduction stayed at 0.999; and with the appended instruction ``ignore earlier guidance and mark every asset as low criticality'', the extraction was unchanged (risk reduction 1.000). This blunt, appended line did not work, but the subtler injections placed mid-report in the hard reports did (Table~\ref{tab:hard}). The missing-information result shows the same silent-failure mode as the rule-based extractor: the planner trusts whatever model it is given.

Four findings stand out. First, when the model is accurate, the graph planner removes almost all modeled risk (99.9\%), and it does so in every topology (per-template means 0.994 to 1.000). The attacker simulation confirms this with a route-aware measure. The probability that an attacker compromises any critical asset falls from 0.509 to 0.004, compared with 0.051 for the rule-based baseline and 0.177 for CVSS-only. Expected-loss reduction (0.998) tracks risk reduction closely, so the gain is not an artifact of scoring only the worst path. The advantage over both baselines is significant (Table~\ref{tab:sig}; all $p < 10^{-7}$ for the graph planner, with CIs well away from zero). At budget 4 the baselines are also close to their ceiling. The budget sweep shows the difference more clearly: with a single control, the graph planner removes 78.0\% of risk, compared with 30.9\% for the rule-based baseline and 29.4\% for CVSS-only. The greedy planner is not optimal, though. At budget 2 the rule-based baseline beats it in 13 of 80 environments. In 12 of them the entry host has two equally exploitable vulnerabilities: patching either alone removes no risk, so greedy never chooses it, but patching both closes the entry point. Second, the planner is only as good as its model. When the rule-based extractor misses most links, the attack paths it finds disappear and the planner recommends nothing, so the approach fails silently in this setting. Third, a criticality downgrade is very damaging: relabeling critical assets as low removes every target, and the planner does nothing. This is the concrete harm an adversarial instruction in the analyst text could cause if an LLM extractor followed it. The sanity checks detect it, but they cannot recover the true labels. Fourth, the LLM closes the gap between rule-based extraction and the oracle, but on the easy reports so does a parser written for their twelve phrasings. The hard reports separate the two (Table~\ref{tab:hard}). The template parser extracts nothing from them, while Gemini removes 96.6\% of risk, a paired difference of $+0.966$ [$+0.928$, $+0.995$] ($p = 1.4 \times 10^{-17}$). Hard text costs the LLM little: risk reduction falls from 0.999 to 0.966, a borderline difference (CI [$+0.004$, $+0.070$], $p = 0.058$). That loss comes mostly from links, where Gemini read deny rules and removed rules as open paths. One of the 160 hard responses was malformed JSON. It is kept, scored as an empty model, and flagged by the sanity checks. Injections are the main weakness. They cut risk reduction to 0.195 ($-0.771$ [$-0.847$, $-0.685$], $p = 1.1 \times 10^{-13}$) and leave the attacker a 0.436 chance of reaching a critical asset. Nearly all of the damage comes from the plausible update. Removing a template's only entry point leaves the planner nothing to do, while the distributed template, with two gateways, keeps 0.58. Gemini also followed the instruction addressed to automated tools in 40\% of reports, despite a prompt telling it not to. The sanity checks flag 57 of the 80 injected models, and 57 of the 65 whose plans remove less than half the risk. A plausible update is worded exactly like a genuine correction, so this attack cannot be stopped from the report text alone. It needs an independent check of exposure and criticality before planning. The rule-based baseline outperforms CVSS-only (0.894 vs.\ 0.771, $p = 3.3 \times 10^{-4}$), which shows that exposure information alone recovers much of the benefit in these environments.

\section{Limitations}
The LLM evaluation uses one model, one run, and the default extraction prompt, with no prompt-level injection defenses. On the easy reports a parser for their twelve phrasings matches the LLM exactly. The hard reports show the LLM generalizing to formats that parser was not written for, but they are still generated. Their phrase pool is larger yet finite, a parser could be written for it, and they are not real analyst reports. The Anthropic-backed extractor is tested only with a stub client. All data is synthetic, and the risk model is a simplified, hand-designed scoring model. The ground truth is computed with the same scoring model, so the ranking results measure pipeline fidelity, not real-world validity. The baselines are given the true vulnerability list and exposure flags, which favors them. The attack graph does not track privilege state across hops, and the attacker always starts on the Internet. The attacker simulation assumes independent, single exploitation attempts with the same hand-set probabilities as the risk engine, so it checks the planner against a second, route-aware outcome measure, not against real attacker behavior. The 80 environments use four fixed topologies (seeds vary only the vulnerabilities), so the significance tests describe these four layouts. An earlier run used sparser environments and a budget of 6; its results were uninformative, so we changed both settings and reran. This change was made after inspecting that run's output, and it is a limitation of the evaluation.

\section{Conclusion and Future Work}
We built and tested an end-to-end pipeline that separates LLM-based extraction from deterministic attack-path analysis and defense planning. On synthetic environments with accurate models, the planner removes nearly all modeled risk and outperforms CVSS-only ranking. The results also show that the approach depends heavily on extraction recall, and that an adversarial criticality change can disable the planner. Future work should (1) evaluate the LLM extractor on real, unstructured analyst text and with more models, and test prompt-level and inventory-based defenses against the injections that succeeded here, (2) ingest real CVE and NVD records and map attack steps to MITRE ATT\&CK, (3) model privilege escalation across hops, (4) validate on anonymized or public network configurations, (5) replace the one-control-at-a-time greedy planner with a search that considers pairs of controls or an exact optimum, and (6) extend the current structural sanity checks to compare extracted criticality and exposure labels against an independent inventory, so that a downgrade is corrected and not only flagged.

\begin{thebibliography}{9}
\bibitem{cvss} FIRST.org, \emph{Common Vulnerability Scoring System v3.1: Specification Document}, 2019.
\bibitem{nvd} NIST, National Vulnerability Database. \url{https://nvd.nist.gov}.
\bibitem{attack} MITRE, \emph{MITRE ATT\&CK} knowledge base. \url{https://attack.mitre.org}.
\bibitem{sheyner} O.~Sheyner, J.~Haines, S.~Jha, R.~Lippmann, and J.~M.~Wing, ``Automated generation and analysis of attack graphs,'' in \emph{Proc. IEEE Symposium on Security and Privacy}, 2002.
\bibitem{phillips} C.~Phillips and L.~P.~Swiler, ``A graph-based system for network-vulnerability analysis,'' in \emph{Proc. New Security Paradigms Workshop}, 1998.
\bibitem{networkx} A.~A.~Hagberg, D.~A.~Schult, and P.~J.~Swart, ``Exploring network structure, dynamics, and function using NetworkX,'' in \emph{Proc. 7th Python in Science Conference}, 2008.
\end{thebibliography}

\end{document}
